\chapter{Конструкторский раздел}

В данном разделе описывается поэтапный процесс взаимодействия с выбранными большими языковыми моделями (БЯМ). Акцент сделан на динамике диалога, количестве итераций и характере возникавших ошибок.

\section{Процесс взаимодействия с моделями}

Для всех моделей использовался единый стартовый запрос (промпт), содержащий полный текст задания и описание варианта. Дальнейшее взаимодействие строилось индивидуально в зависимости от качества первичного ответа.

\subsection{Взаимодействие с Gemini 3 Pro Preview}
Взаимодействие с моделью Gemini показало наивысшую эффективность, процесс занял минимальное количество шагов.

\begin{enumerate}
    \item \textbf{Первичная генерация:} модель с первой попытки выдала синтаксически корректный код и составила корректные регулярные выражения;
    \item \textbf{Анализ недостатков:} в исходном решении поиск производился построчно. Это создавало риск пропуска ошибок, если словосочетание (например, <<рисунок 1>>) было разорвано переносом строки в самом PDF-документе;
    \item \textbf{Внесение корректировок:} был подан один дополнительный уточняющий запрос с требованием изменить логику поиска с построчной на проверку полного текста страницы целиком;
    \item \textbf{Результат:} модель успешно переписала функцию, что позволило получить итоговое рабочее решение.
\end{enumerate}

\subsection{Взаимодействие с Qwen 3 Coder (Local)}
Работа с локальной моделью потребовала большего участия оператора, что объясняется меньшим объёмом параметров модели по сравнению с облачными аналогами.

\begin{enumerate}
    \item \textbf{Первичная генерация:} модель сгенерировала код, который запускался, но содержал логическую инверсию задачи. Вместо поиска \textit{отсутствия} неразрывного пробела (ошибки), регулярные выражения искали его \textit{наличие} (правильный вариант);
    \item \textbf{Итеративное уточнение (логика):} потребовался дополнительный промпт для объяснения концепции <<поиска ошибки>>;
    \item \textbf{Техническая доработка:} после исправления логики регулярных выражений потребовалось дополнительное указание на необходимость поиска ``от'' и ``до'' по отдельности, а не в составе единой конструкции. Также потребовалось указать на необходимость игнорировать регистр, исправить логику поиска с построчного, на поиск по всему тексту;
    \item \textbf{результат:} после трёх итераций уточнений был получен корректный скрипт.
\end{enumerate}

\subsection{Взаимодействие с YandexGPT 5.1 Pro}
Взаимодействие с данной моделью выявило специфическую проблему <<инерции мышления>>: модель успешно исправляла логику регулярных выражений, но игнорировала структурные ошибки в коде Python.

\begin{enumerate}
    \item \textbf{Первичная генерация:} модель сгенерировала код с грубой ошибкой управления ресурсами \texttt{(Context Management)}. Файл открывался через \texttt{with open(...)}, создавался объект \texttt{PdfReader}, после чего блок закрывался, а попытка чтения страниц происходила уже вне контекста, что приводило к ошибке \texttt{I/O operation on closed file};
    \item \textbf{Логическая ошибка валидации:} изначально модель неверно интерпретировала условие задачи, считая успехом наличие неразрывного пробела, тогда как требовалось найти места его \textit{отсутствия}. Это было исправлено после первого замечания;
    \item \textbf{Циклическое игнорирование ошибки:} несмотря на три последовательных прямых указания на то, что переменные \texttt{reader} используются вне блока \texttt{with}, модель раз за разом генерировала <<исправленный>> код, сохраняя ошибочную структуру отступов. Исправления касались только регулярных выражений (добавление \texttt{re.IGNORECASE}, разделение паттернов);
    \item \textbf{Результат:} корректный код, в котором цикл обработки страниц находится внутри блока открытия файла, был получен только после четвёртого запроса, после жёсткого указания на то, что ошибка не была исправлена.
\end{enumerate}

\section*{Выводы}
Сравнительный анализ процесса разработки показал:
\begin{itemize}
    \item \textbf{Gemini 3 Pro} продемонстрировала способность к <<автономной разработке>>: минимум правок, глубокое понимание контекста задачи;
    \item \textbf{Qwen 3 Coder} эффективна в роли <<ассистента>>, но требует чёткого технического надзора и декомпозиции задачи на простые логические шаги;
    \item \textbf{YandexGPT 5.1} на текущем этапе требует внимательной валидации каждой строки кода, так как склонна игнорировать инструкции по исправлению собственных ошибок.
\end{itemize}

\clearpage